% Textos fijos, espa\~nol (versi\'on MSX2).
\usepackage[spanish]{babel}
\renewcommand{\portWork}{\portWorkES}
\newcommand{\tManual}{Manual de usuario}
\newcommand{\tUserManual}{MANUAL DE USUARIO}
\newcommand{\tForV}{para MSX2}
\newcommand{\tForTheV}{para MSX2}
\newcommand{\tCoverLine}{El juego del Spectrum, en color en tu MSX2}
\newcommand{\tOriginalBy}{Un juego original de}
\newcommand{\tPortBy}{Versi\'on para MSX2 de}
\newcommand{\tHomage}{Esta es una versi\'on no oficial hecha para el estudio, la conservaci\'on y como homenaje.}
\newcommand{\tWhatYouNeed}{Qu\'e necesitas}
\newcommand{\tItem}{Elemento}
\newcommand{\tRequirement}{Requisito}
\newcommand{\tComputer}{Ordenador}
\newcommand{\tComputerText}{\textbf{Cualquier MSX2} (o MSX2+ o turbo R) con una ranura de cartucho libre. Funciona a 50 y a 60~Hz y el juego va a la misma velocidad en los dos.}
\newcommand{\tMemory}{Memoria}
\newcommand{\tMemoryText}{\textbf{\ramNeeded{} de RAM} y \textbf{128~KB de VRAM} (lo normal en un MSX2).}
\newcommand{\tVideo}{V\'ideo}
\newcommand{\tVideoText}{El chip de v\'ideo del MSX2 (V9938 o V9958), en SCREEN~5. No hace falta ning\'un cartucho de v\'ideo.}
\newcommand{\tSoundRow}{Sonido}
\newcommand{\tSoundText}{El PSG del MSX: la m\'usica (t\'itulo y partida) y los efectos de sonido.}
\newcommand{\tGameRow}{Juego}
\newcommand{\tGameText}{El \textbf{cartucho de \gameTitle{} para MSX2}: una \romKind{} (el fichero \texttt{\romFile}, grabado en un cartucho flash o equivalente).}
\newcommand{\tMonitor}{Monitor}
\newcommand{\tMonitorText}{El monitor o televisor del MSX.}
\newcommand{\tControlsRow}{Controles}
\newcommand{\tControlsText}{Cursores y barra espaciadora, o un joystick en el puerto 1 o en el 2.}
\newcommand{\tNoTape}{\textbf{Sin cinta, sin disco, sin contrase\~na.} Todo el juego est\'a en el cartucho.}
\newcommand{\tSettingUp}{Puesta en marcha}
\newcommand{\tSetOff}{\textbf{Apaga el MSX.} No insertes ni retires nunca un cartucho con el equipo encendido.}
\newcommand{\tSetGame}{Inserta el \textbf{cartucho de \gameTitle{}} en una ranura de cartucho.}
\newcommand{\tSetMonitor}{Conecta el monitor o televisor como siempre.}
\newcommand{\tSetJoy}{Conecta un joystick al puerto 1 o al 2 si quieres usarlo.}
\newcommand{\tSetOn}{Enciende el MSX. Tras el arranque del ordenador aparece el logotipo de DiHalt Studio.}
\newcommand{\tBlackScreen}{\textbf{Si la pantalla se queda en negro} comprueba que el cartucho est\'a bien insertado y que el ordenador es un MSX2 con \ramNeeded{} de RAM. En un emulador como openMSX, carga la ROM con el mapper \textbf{ASCII16} (opci\'on \texttt{-romtype ASCII16}).}
\newcommand{\tFlashCart}{\textbf{Con un MegaFlashROM SCC+ SD:} copia en la ra\'iz de la tarjeta SD el contenido de la carpeta \texttt{dist} (el juego como \texttt{ANDYCAPP.ROM}, el programa \texttt{OPFXSD}, los ficheros para arrancar desde la tarjeta y un \texttt{AUTOEXEC.BAT}) y enciende el MSX: el juego se graba en el cartucho en unos segundos y empieza. Desde entonces el cartucho arranca \gameTitle{} cada vez que enciendes. Para arrancar desde la tarjeta mant\'en pulsada \key{ESC} al encender (tambi\'en la primera vez, si el cartucho ya ten\'ia otro juego grabado); \texttt{OPFXSD /C} borra el juego del cartucho.}
\newcommand{\tPowerOnTitle}{Del encendido al juego}
\newcommand{\tPowerOnCaption}{Izquierda: la portada. Derecha: la pantalla de t\'itulo.}
\newcommand{\tControlsTitle}{Controles}
\newcommand{\tCreditsTitle}{9. Cr\'editos y derechos}
\newcommand{\tIsAGameBy}{es un juego de}
\newcommand{\tRightsText}{Programado por Blitter Animations; m\'usica de Jason C.~Brooke. Los personajes de Andy Capp son creaci\'on de Reg Smythe y propiedad de Mirror Group Newspapers. El juego, sus gr\'aficos, m\'usica y nombre pertenecen a sus titulares. Esta es una versi\'on no oficial hecha para el estudio, la conservaci\'on y como homenaje; no tiene relaci\'on con ellos.}
\newcommand{\tVersionWord}{Versi\'on MSX2}
\newcommand{\tMsxTrademark}{MSX es una marca de MSX Licensing Corporation. El logotipo MSX se muestra solo para indicar en qu\'e ordenadores funciona el juego}
\newcommand{\tSourceNote}{C\'odigo fuente disponible en GitHub. \textbf{\textexclamdown Mejora el juego si te apetece!}}
\newcommand{\tAppendix}{Ap\'endice. Etiqueta del cartucho}
\newcommand{\tStickerText}{Una etiqueta para el cartucho, lista para imprimir. Al 100\,\% mide unos 16\,cm de ancho; esc\'alala en tu programa de impresi\'on al \'area de etiqueta de tu cartucho (la imagen tiene proporci\'on 3:2).}
\newcommand{\tAction}{Acci\'on}
\newcommand{\tKeyboard}{Teclado}
\newcommand{\tJoystick}{Joystick}
\newcommand{\tNo}{N.\textordmasculine}
\newcommand{\tElement}{Elemento}
\newcommand{\tTellsYou}{Qu\'e te indica}
\newcommand{\tTips}{Consejos.}
